\documentclass[11pt,a4paper]{ctexart}
\usepackage[teacher]{ipara}
\tcbuselibrary{breakable}
\usepackage{multicol}
\raggedbottom


\pagestyle{fancy}
\fancyhf{}
\lhead{\small 高一/高二拔高专题：立体几何折叠与折痕问题}
\rhead{\small 教师版（不变量、位置关系、距离、最值与轨迹）}
\cfoot{\small 第 \thepage\ 页\quad 共 \zpageref{LastPage} 页}
\renewcommand{\headrulewidth}{0.4pt}
\renewcommand{\footrulewidth}{0pt}

% 自定义卡片与结构环境

\renewcommand{\teacherproblem}[2]{%
  \teachquestionheader{例题 #1\quad #2}
}

% 引入外部公共题目宏定义以备用（我们已经录入了公共题）
% =====================================================
% 题型：立体几何解答题 / 折叠与最值问题 (q_solid_folding_rectangle.tex)
% =====================================================
% =====================================================
% 难度：★★★ (难题)
% =====================================================
\newcommand{\qSolidFoldingRectangleStem}{%
  如图，在矩形 $ABCD$ 中，$AB=1$，$BC=\sqrt{3}$，$M$ 是线段 $AD$ 上的一动点，将 $\triangle ABM$ 沿着 $BM$ 折起，使点 $A$ 到达点 $A'$ 的位置，满足点 $A' \notin$ 平面 $BCDM$，且点 $A'$ 在平面 $BCDM$ 内的射影 $E$ 落在线段 $BC$ 上。
  \begin{enumerate}[label=(\arabic*), leftmargin=2em, itemsep=0.1em]
    \item 当点 $M$ 与端点 $D$ 重合时，证明：$A'B \perp$ 平面 $A'CD$；
    \item 当 $AM=\dfrac{3}{2}$ 时，求二面角 $A'-BM-C$ 的余弦值；
    \item 设直线 $CD$ 与平面 $A'BM$ 所成的角为 $\alpha$，二面角 $A'-BM-C$ 的平面角为 $\beta$，求 $\sin^2\alpha \cos\beta$ 的最大值。
  \end{enumerate}%
}

\newcommand{\qSolidFoldingRectangleFigure}[1][1.0]{%
  \begin{tikzpicture}[scale=#1, >=stealth, line join=round, line cap=round]
    % 左图：折叠前
    \begin{scope}[xshift=0cm, yshift=0cm]
      % 定义顶点
      \coordinate (B) at (0, 0);
      \coordinate (C) at (2.6, 0);
      \coordinate (D) at (2.6, 1.5);
      \coordinate (A) at (0, 1.5);
      \coordinate (M) at (1.8, 1.5);
      \coordinate (E) at (1.25, 0);
      \coordinate (F) at (0.74, 0.62);
      
      % 绘制矩形和折痕
      \draw[thick] (A) -- (B) -- (C) -- (D) -- cycle;
      \draw[thick, dashed, blue] (B) -- (M);
      \draw[thick, dashed, red] (A) -- (E);
      
      % 标记点
      \fill (A) circle (1.2pt) node[above left] {\small $A$};
      \fill (B) circle (1.2pt) node[below left] {\small $B$};
      \fill (C) circle (1.2pt) node[below right] {\small $C$};
      \fill (D) circle (1.2pt) node[above right] {\small $D$};
      \fill (M) circle (1.2pt) node[above] {\small $M$};
      \fill (E) circle (1.2pt) node[below] {\small $E$};
      \fill (F) circle (1.2pt) node[above left] {\small $F$};
      
      % 直角标记
      \draw[gray, thin] ($(F)!0.1!(A)$) -- ($($(F)!0.1!(A)$) + ($(F)!0.1!(M)$) - (F)$) -- ($(F)!0.1!(M)$);
      
      \node[below] at (1.3, -0.2) {\small 折叠前};
    \end{scope}
    
    % 折叠箭头
    \draw[->, very thick, gray] (3.0, 0.75) -- (4.0, 0.75) node[midway, above] {\small 折叠};
    
    % 右图：折叠后
    \begin{scope}[xshift=4.5cm, yshift=0cm]
      % 定义顶点 (精确计算的投影三维坐标)
      \coordinate (B) at (0, 0);
      \coordinate (C) at (2.6, 0);
      \coordinate (D) at (3.2, 0.9);
      \coordinate (M) at (2.55, 0.9);
      \coordinate (E) at (1.15, 0);
      \coordinate (A1) at (1.15, 0.96);
      \coordinate (F) at (1.02, 0.36); % BM 上的垂足 F
      
      % 绘制底面边界与棱
      \draw[thick] (B) -- (C) -- (D);
      \draw[thick, dashed, gray!60] (D) -- (M);
      \draw[thick, dashed, gray!60] (M) -- (B);
      
      % 绘制折起后的面 A'BM 和 A'BC
      \draw[thick, blue] (A1) -- (B);
      \draw[thick, blue] (A1) -- (M);
      \draw[thick] (A1) -- (C);
      \draw[thick, dashed, gray!60] (A1) -- (D);
      
      % 绘制辅助线
      \draw[thick, dashed, red] (A1) -- (E);
      \draw[thick, dashed, red] (E) -- (F);
      \draw[thick, dashed, red] (A1) -- (F);
      
      % 标记点
      \fill (B) circle (1.2pt) node[below left] {\small $B$};
      \fill (C) circle (1.2pt) node[below right] {\small $C$};
      \fill (D) circle (1.2pt) node[above right] {\small $D$};
      \fill (M) circle (1.2pt) node[above right] {\small $M$};
      \fill (E) circle (1.2pt) node[below] {\small $E$};
      \fill (A1) circle (1.2pt) node[above] {\small $A'$};
      \fill (F) circle (1.2pt) node[left] {\small $F$};
      
      % 直角标记
      % A'E \perp base
      \draw[red, thin] ($(E)!0.08!(A1)$) -- ($($(E)!0.08!(A1)$) + ($(E)!0.08!(B)$) - (E)$) -- ($(E)!0.08!(B)$);
      % A'F \perp BM
      \draw[red, thin] ($(F)!0.08!(A1)$) -- ($($(F)!0.08!(A1)$) + ($(F)!0.08!(M)$) - (F)$) -- ($(F)!0.08!(M)$);
      % EF \perp BM
      \draw[red, thin] ($(F)!0.08!(E)$) -- ($($(F)!0.08!(E)$) + ($(F)!0.08!(M)$) - (F)$) -- ($(F)!0.08!(M)$);
      
      \node[below] at (1.3, -0.2) {\small 折叠后};
    \end{scope}
  \end{tikzpicture}%
}

\newcommand{\qSolidFoldingRectangleAnswer}{%
  (1) 证明见解析；(2) 二面角 $A'-BM-C$ 的余弦值为 $\dfrac{4}{9}$；(3) $\sin^2\alpha \cos\beta$ 的最大值为 $\dfrac{1}{4}$.%
}

\newcommand{\qSolidFoldingRectangleSolution}{%
  【解题思路】\\
  本题为折叠几何综合最值难题。解题的关键在于抓住折叠前后的“距离与角度不变性”，通过代数建模构建起动态线段的长度关系网络：
  \begin{enumerate}[label=\arabic*.]
    \item \textbf{建立核心长度模型}：设 $AM=x$。通过 $A'$ 在底面的射影 $E$ 落在 $BC$ 上这一临界条件，结合折叠前后 $A'B=AB=1, A'M=AM=x$ 及 $A'E \perp$ 平面 $BCDM$，在底面利用 Rt$\triangle A'EB$ 和 Rt$\triangle A'EM$ 构建方程，推导得出 $BE = \dfrac{1}{x}$，此式为全题的计算发动机。
    \item \textbf{几何垂直证明}：当 $M=D$ 时，长度参数 $x=\sqrt{3}$ 确定。利用勾股定理逆定理在 $\triangle A'BC$ 中求证 $A'B \perp A'C$，再利用线面垂直定理证明 $CD \perp$ 平面 $A'BC$ 从而得到 $CD \perp A'B$，最终证明 $A'B \perp$ 平面 $A'CD$。
    \item \textbf{二面角计算}：利用三垂线定理判定二面角 $A'-BM-C$ 的平面角为 $\angle A'FE$。通过底面三角形面积法计算出 $EF$，再结合折叠前后 $A'F=AF$（即 $A$ 到折痕 $BM$ 的距离），在直角三角形 $\triangle A'EF$ 中直接求出 $\cos\beta = \dfrac{EF}{A'F} = \dfrac{1}{x^2}$。
    \item \textbf{最值求解与三正弦定理模型识别}：注意到直线 $CD$ 在底面内，平面 $A'BM$ 是底面绕 $BM$ 折起形成的。根据\textbf{三正弦定理（最大角定理）}，直线与折起平面所成的角 $\alpha$、折起二面角 $\beta$ 以及直线与折痕的夹角 $\theta$ 满足：$\sin\alpha = \sin\beta \sin\theta$。由此将线面角转化为关于 $x$ 的代数式，化简后应用二次函数求最值。
  \end{enumerate}

  【解析计算】\\
  \noindent\textbf{核心关系推导：}\\
  设 $AM=x$，$A'$ 在底面上的射影为 $E \in BC$，则 $A'E \perp$ 平面 $BCDM$。\\
  因为折叠是绕 $BM$ 进行的刚体旋转，折叠前后距离保持不变，故有：
  \[ A'B = AB = 1,\quad A'M = AM = x \]
  由于 $A'E \perp$ 平面 $BCDM$，且 $B, C, E \subset$ 平面 $BCDM$，所以 $A'E \perp BC$。\\
  在 Rt$\triangle A'EB$ 中，有：
  \[ A'E^2 = A'B^2 - BE^2 = 1 - BE^2 \quad \text{—— 方程 (1)} \]
  在底面矩形中，过 $M$ 作 $MM_0 \perp BC$ 于 $M_0$，由矩形性质可知 $BM_0 = AM = x, MM_0 = AB = 1$。\\
  由于 $E \in BC$，在 Rt$\triangle EM_0M$ 中，有：
  \[ EM^2 = EM_0^2 + MM_0^2 = (BE - x)^2 + 1 \]
  在 Rt$\triangle A'EM$ 中，因为 $A'E \perp$ 平面 $BCDM$，所以 $A'E \perp EM$。则有：
  \[ A'M^2 = EM^2 + A'E^2 \implies x^2 = (BE - x)^2 + 1 + A'E^2 \quad \text{—— 方程 (2)} \]
  将方程 (1) 代入方程 (2)，消去 $A'E^2$ 得：
  \[ x^2 = (BE - x)^2 + 1 + (1 - BE^2) \]
  展开整理得：
  \[ x^2 = BE^2 - 2x \cdot BE + x^2 + 2 - BE^2 \implies 2x \cdot BE = 2 \implies \boxed{BE = \frac{1}{x}} \]
  代入方程 (1) 中，可得 $A'$ 的投影高度为：
  \[ \boxed{A'E = \sqrt{1 - \frac{1}{x^2}}} \]
  要使折起后的几何体存在，高 $A'E$ 必须为实数且不为 $0$，即 $1 - \frac{1}{x^2} > 0 \implies x > 1$。\\
  又因为点 $M$ 在线段 $AD$ 上，且 $AD = \sqrt{3}$，所以 $x \le \sqrt{3}$。\\
  故动点参数 $x$ 的取值范围为 $\boxed{x \in [1, \sqrt{3}]}$。

  \noindent\textbf{（1）证明：当 $M=D$ 时，$A'B \perp$ 平面 $A'CD$}\\
  当点 $M$ 与端点 $D$ 重合时，其长度参数 $x = AM = AD = BC = \sqrt{3}$。\\
  此时可得：
  \[ BE = \frac{1}{\sqrt{3}},\quad CE = BC - BE = \sqrt{3} - \frac{1}{\sqrt{3}} = \frac{2}{\sqrt{3}} \]
  高度的平方为：
  \[ A'E^2 = 1 - \frac{1}{x^2} = 1 - \frac{1}{3} = \frac{2}{3} \]
  在 Rt$\triangle A'EC$ 中，有：
  \[ A'C^2 = A'E^2 + CE^2 = \frac{2}{3} + \left(\frac{2}{\sqrt{3}}\right)^2 = \frac{2}{3} + \frac{4}{3} = 2 \]
  在 $\triangle A'BC$ 中，因为 $A'B = 1, BC = \sqrt{3}$，且 $A'C = \sqrt{2}$。\\
  易见：
  \[ A'B^2 + A'C^2 = 1^2 + (\sqrt{2})^2 = 3 = BC^2 \]
  根据勾股定理逆定理，可知 $\boxed{A'B \perp A'C}$。\\
  在底面矩形中，$CD \perp BC$。又因为 $A'E \perp$ 平面 $BCDM$，且 $CD \subset$ 平面 $BCDM$，所以 $CD \perp A'E$。\\
  由于 $BC \cap A'E = E$，所以 $CD \perp$ 平面 $A'BC$。\\
  因为 $A'B \subset$ 平面 $A'BC$，所以有 $\boxed{CD \perp A'B}$。\\
  由于 $A'C \cap CD = C$，且 $A'C, CD \subset$ 平面 $A'CD$，\\
  综上所述，$A'B$ 垂直于平面 $A'CD$ 内的两条相交直线，故有：
  \[ \boxed{A'B \perp \text{平面 } A'CD} \]

  \noindent\textbf{（2）求解二面角 $A'-BM-C$ 的平面角：}\\
  在底面内，过点 $E$ 作 $EF \perp BM$ 于点 $F$，连接 $A'F$。\\
  因为 $A'E \perp$ 平面 $BCDM$，且 $EF \perp BM$，根据三垂线定理可知，$\boxed{A'F \perp BM}$。\\
  因此，$\angle A'FE$ 即为二面角 $A'-BM-C$ 的平面角 $\beta$。\\
  在直角三角形 $\triangle ABM$ 中，斜边为 $BM = \sqrt{AM^2 + AB^2} = \sqrt{x^2 + 1}$。\\
  在底面中，三角形 $\triangle BEM$ 的面积可以表示为：
  \[ S_{\triangle BEM} = \frac{1}{2} \cdot BE \cdot AB = \frac{1}{2} \cdot \frac{1}{x} \cdot 1 = \frac{1}{2x} \]
  同时，以 $BM$ 为底，高为 $EF$，其面积也可以表示为：
  \[ S_{\triangle BEM} = \frac{1}{2} \cdot BM \cdot EF = \frac{1}{2} \cdot \sqrt{x^2+1} \cdot EF \]
  联立面积公式，求得 $EF$ 的长度为：
  \[ EF = \frac{1}{x\sqrt{x^2+1}} \]
  折叠前，点 $A$ 到折痕 $BM$ 的距离等于折起后点 $A'$ 到折痕 $BM$ 的距离（即 $A'F$）。\\
  在直角三角形 $\triangle ABM$ 中，由等面积法有：
  \[ S_{\triangle ABM} = \frac{1}{2} \cdot AB \cdot AM = \frac{x}{2} = \frac{1}{2} \cdot BM \cdot A'F \implies A'F = \frac{x}{\sqrt{x^2+1}} \]
  在 Rt$\triangle A'EF$ 中，二面角 $\beta = \angle A'FE$ 的余弦值为：
  \[ \cos\beta = \frac{EF}{A'F} = \frac{\frac{1}{x\sqrt{x^2+1}}}{\frac{x}{\sqrt{x^2+1}}} = \boxed{\frac{1}{x^2}} \]
  当 $AM = \frac{3}{2}$ 即 $x = \frac{3}{2}$ 时，代入可得：
  \[ \cos\beta = \frac{1}{\left(\frac{3}{2}\right)^2} = \frac{4}{9} \]
  故所求二面角 $A'-BM-C$ 的余弦值为 $\boxed{\dfrac{4}{9}}$。

  \noindent\textbf{（3）求 $\sin^2\alpha \cos\beta$ 的最大值：}\\
  在底面矩形中，直线 $CD \parallel AB$，所以直线 $CD$ 与折痕 $BM$ 所成的角 $\theta$ 等于直线 $AB$ 与 $BM$ 所成的角。\\
  在直角三角形 $\triangle ABM$ 中，有：
  \[ \sin\theta = \frac{AM}{BM} = \frac{x}{\sqrt{x^2+1}} \]
  直线 $CD$ 位于底面内，平面 $A'BM$ 是由底面的一部分绕 $BM$ 折起形成的，且这两个平面的夹角为 $\beta$。\\
  根据\textbf{三正弦定理（最大角定理）}，直线 $CD$ 与平面 $A'BM$ 所成的线面角 $\alpha$ 满足关系式：
  \[ \sin\alpha = \sin\beta \sin\theta \]
  由（2）中求得 $\cos\beta = \frac{1}{x^2}$，由于 $x \in [1, \sqrt{3}]$，故有：
  \[ \sin\beta = \sqrt{1 - \cos^2\beta} = \sqrt{1 - \frac{1}{x^4}} \]
  代入可求得：
  \[ \sin^2\alpha = \sin^2\beta \sin^2\theta = \left(1 - \frac{1}{x^4}\right) \cdot \frac{x^2}{x^2+1} = \frac{x^4-1}{x^4} \cdot \frac{x^2}{x^2+1} \]
  利用平方差公式进行化简：
  \[ \sin^2\alpha = \frac{(x^2-1)(x^2+1)}{x^4} \cdot \frac{x^2}{x^2+1} = \frac{x^2-1}{x^2} \]
  由此可得目标代数式为：
  \[ \sin^2\alpha \cos\beta = \frac{x^2-1}{x^2} \cdot \frac{1}{x^2} = \frac{x^2-1}{x^4} \]
  令 $t = x^2$。因为 $x \in [1, \sqrt{3}]$，所以 $t \in [1, 3]$。\\
  则代数式化为：
  \[ f(t) = \frac{t-1}{t^2} = \frac{1}{t} - \frac{1}{t^2} \]
  进一步令 $u = \frac{1}{t}$。因为 $t \in [1, 3]$，所以 $u \in \left[\frac{1}{3}, 1\right]$。\\
  目标二次函数为：
  \[ g(u) = u - u^2 = -\left(u - \frac{1}{2}\right)^2 + \frac{1}{4} \]
  因为对称轴 $u = \frac{1}{2} \in \left[\frac{1}{3}, 1\right]$，所以当 $u = \frac{1}{2}$（即 $t = 2 \implies x = \sqrt{2} \in [1, \sqrt{3}]$）时，函数取得最大值，最大值为：
  \[ g\left(\frac{1}{2}\right) = \frac{1}{4} \]
  故所求 $\sin^2\alpha \cos\beta$ 的最大值为 $\boxed{\dfrac{1}{4}}$。
}

\newcommand{\qSolidFoldingRectangleDiagnosis}{%
  用于课堂精准提优训练。本题是一道非常经典的矩形折叠最值难题。重点考查折叠前后“长度与角度的守恒性”以及将动态几何特征转化为代数方程的能力。在解析中应引导学生识别出“三正弦定理”在求线面角最值中的应用，免去在折叠后的复杂空间体系中寻找线面角投影的繁琐作图过程。
}

\newcommand{\qSolidFoldingRectangleTeachingNotes}{%
  \item \textbf{教学核心与模型识别提示}：
    \begin{itemize}
      \item \textbf{折叠动力引擎}：引导学生先建立折起高 $A'E = \sqrt{1 - \frac{1}{x^2}}$ 与投影点位置 $BE = \frac{1}{x}$ 的代数关系，这是折叠模型中利用垂直平分性质建立等式的一个非常典型的“保距消元”手法，务必让学生彻底掌握这一推导。
      \item \textbf{三正弦定理的妙用}：第（3）问中，如果按照常规思路去寻找直线 $CD$ 在平面 $A'BM$ 上的投影来求线面角 $\alpha$，空间想象难度极大且极易出错。这里直接把底面看作平面 $\beta$，平面 $A'BM$ 看作平面 $\alpha$，直线 $CD$ 位于底面内，利用\textbf{三正弦定理} $\sin\alpha = \sin\beta \sin\theta$（线面角正弦 = 二面角正弦 $\times$ 面内线与交线夹角正弦），直接跳过了复杂的空间垂足寻找过程，是极佳的模型识别实战案例。
    \end{itemize}
}

% 难度：★★★
% 知识点：立体几何、折叠问题、线面垂直、外接球、二面角最值

\newcommand{\qSolidFoldingPentagonMinAngleStem}{
在凸五边形 $ABCDE$ 中，$\angle EAC = \angle DCA = \frac{\pi}{2}$，$AE = 2CD = 4$，$\triangle ABC$ 是边长为 $2\sqrt{3}$ 的正三角形。点 $F, H$ 分别为 $ED, AD$ 的中点。将 $\triangle ABC$ 沿 $AC$ 翻折，点 $B$ 到空间点 $P$。
\begin{enumerate}[label=(\arabic*)]
  \item 若 $PH \perp$ 平面 $ACDE$，求：
  \begin{enumerate}[label=\textcircled{\arabic*}]
    \item 直线 $PA$ 与 $CF$ 所成的角；
    \item 三棱锥 $P-ACD$ 的外接球为球 $O$，$M$ 为球面上一动点，求 $EM$ 的取值范围。
  \end{enumerate}
  \item 在翻折过程中，设平面 $PCD$ 与平面 $PAE$ 的交线为 $l$，求二面角 $A-l-C$ 的最小值。
\end{enumerate}
}

\newcommand{\qSolidFoldingPentagonMinAngleFigure}[1][1]{
\begin{tikzpicture}[scale=#1, >=stealth, line join=round, line cap=round]
  % 2D 平面图
  \begin{scope}[xshift=0cm, yshift=0cm]
    \coordinate (A) at (0,0);
    \coordinate (C) at (3.464,0);
    \coordinate (E) at (0,4);
    \coordinate (D) at (3.464,2);
    \coordinate (B) at (1.732,-3);
    \coordinate (H) at (1.732,1);
    \coordinate (F) at (1.732,3);
    
    \draw[thick] (B) -- (C) -- (D) -- (E) -- (A) -- cycle;
    \draw[thick] (A) -- (C);
    \draw[thin, dashed] (A) -- (D);
    \draw[thin, dashed] (C) -- (F);
    
    \fill (A) circle (1.2pt) node[left] {$A$};
    \fill (C) circle (1.2pt) node[right] {$C$};
    \fill (E) circle (1.2pt) node[left] {$E$};
    \fill (D) circle (1.2pt) node[right] {$D$};
    \fill (B) circle (1.2pt) node[below] {$B$};
    \fill (H) circle (1.2pt) node[below right] {$H$};
    \fill (F) circle (1.2pt) node[above right] {$F$};
    \node[below] at (1.732,-3.5) {\small 翻折前平面图};
  \end{scope}

  % 3D 折叠图
  \begin{scope}[xshift=5.5cm, yshift=-1.0cm]
    \coordinate (A) at (0,0);
    \coordinate (C) at (3.5,0);
    \coordinate (E) at (1.2,1.6);
    \coordinate (D) at (4.1,0.8);
    
    \coordinate (H) at (2.05,0.4); % AD 中点
    \coordinate (F) at (2.65,1.2); % ED 中点
    \coordinate (P) at (2.05, 3.5); % P在 H 正上方
    
    % 实线
    \draw[thick] (A) -- (C) -- (D);
    \draw[thick] (P) -- (A);
    \draw[thick] (P) -- (C);
    \draw[thick] (P) -- (D);
    
    % 虚线
    \draw[thick, dashed] (A) -- (E) -- (D);
    \draw[thick, dashed] (P) -- (E);
    \draw[thin, dashed] (A) -- (D);
    \draw[thin, dashed] (P) -- (H);
    \draw[thin, dashed] (C) -- (F);
    
    % 垂直标记
    \draw[thin] ($(H)!0.08!(P)$) -- ($($(H)!0.08!(P)$) + ($(H)!0.1!(A)$) - (H)$) -- ($(H)!0.1!(A)$);
    
    \fill (A) circle (1.2pt) node[below left] {$A$};
    \fill (C) circle (1.2pt) node[below] {$C$};
    \fill (E) circle (1.2pt) node[left] {$E$};
    \fill (D) circle (1.2pt) node[right] {$D$};
    \fill (P) circle (1.2pt) node[above] {$P$};
    \fill (H) circle (1.2pt) node[below] {$H$};
    \fill (F) circle (1.2pt) node[above right] {$F$};
    \node[below] at (2,-0.5) {\small 翻折后三维图};
  \end{scope}
\end{tikzpicture}
}

\newcommand{\qSolidFoldingPentagonMinAngleAnswer}{
(1) \textcircled{1} $\frac{\pi}{2}$ \quad \textcircled{2} $[\sqrt{2}, 4\sqrt{2}]$ \quad (2) $\frac{\pi}{3}$
}

\newcommand{\qSolidFoldingPentagonMinAngleSolution}{
\textbf{（1）} \textcircled{1} \textbf{求 $PA$ 与 $CF$ 所成角：}\\
因为 $AE \perp AC$，$CD \perp AC$，所以 $AE \parallel CD$。\\
又 $AC = 2\sqrt{3}, AE = 4, CD = 2$，所以 $ACDE$ 是直角梯形。\\
设 $G$ 为 $AC$ 的中点，则 $AG = CG = \sqrt{3}$。\\
因为 $F$ 是 $ED$ 的中点，根据直角梯形中位线性质，$F$ 在 $AC$ 上的射影为 $G$，且 $GF = \frac{AE+CD}{2} = 3$。\\
因为 $H$ 是 $AD$ 的中点，$H$ 在 $AC$ 上的射影也是 $G$，且 $GH = \frac{CD}{2} = 1$。\\
在平面 $ACDE$ 内，$\overrightarrow{AH} = \overrightarrow{AG} + \overrightarrow{GH}$，$\overrightarrow{CF} = \overrightarrow{CG} + \overrightarrow{GF}$。\\
由于 $AG \perp GH$，$CG \perp GF$，且 $\overrightarrow{AG}$ 与 $\overrightarrow{CG}$ 方向相反，可得：
\[
\overrightarrow{AH} \cdot \overrightarrow{CF} = \overrightarrow{GH} \cdot \overrightarrow{GF} - |\overrightarrow{AG}| \cdot |\overrightarrow{CG}| = 1 \times 3 - \sqrt{3} \times \sqrt{3} = 0
\]
所以 $AH \perp CF$。\\
又因为 $PH \perp$ 平面 $ACDE$，$CF \subset$ 平面 $ACDE$，所以 $PH \perp CF$。\\
因为 $AH \cap PH = H$，所以 $CF \perp$ 平面 $APH$。\\
因为 $PA \subset$ 平面 $APH$，所以 $CF \perp PA$。\\
故 $PA$ 与 $CF$ 所成角为 $\frac{\pi}{2}$。

\textcircled{2} \textbf{求 $EM$ 的取值范围：}\\
因为 $AC \perp CD$，$AC = 2\sqrt{3}, CD = 2$，所以 $AD = \sqrt{12+4} = 4$。\\
所以 $\triangle ACD$ 是直角三角形，其外接圆圆心为斜边 $AD$ 的中点 $H$。\\
因为 $PH \perp$ 平面 $ACDE$，所以三棱锥 $P-ACD$ 的外接球球心 $O$ 必在直线 $PH$ 上。\\
设 $OH = t$。在 $\triangle PAH$ 中，$PA = 2\sqrt{3}, AH = 2$，则 $PH = \sqrt{PA^2 - AH^2} = 2\sqrt{2}$。\\
由球心性质 $OA = OP$，有：
\[
AH^2 + OH^2 = (PH - OH)^2 \implies 4 + t^2 = (2\sqrt{2} - t)^2
\]
解得 $t = \frac{1}{\sqrt{2}}$。即球半径 $R = OA = \sqrt{4 + \frac{1}{2}} = \frac{3}{\sqrt{2}}$。\\
在平面 $ACDE$ 内，建立以 $A$ 为原点，$AC$ 为 $x$ 轴的直角坐标系，则 $A(0,0)$，$E(0,4)$，$D(2\sqrt{3}, 2)$，所以 $H(\sqrt{3}, 1)$。\\
所以 $EH^2 = (\sqrt{3} - 0)^2 + (1 - 4)^2 = 12$。\\
在空间中，$OH \perp$ 平面 $ACDE$，所以：
\[
EO^2 = EH^2 + OH^2 = 12 + \frac{1}{2} = \frac{25}{2} \implies EO = \frac{5}{\sqrt{2}}
\]
点 $M$ 在球面上运动，故 $EM$ 的最大值与最小值分别为：
\[
EM_{\max} = EO + R = \frac{5}{\sqrt{2}} + \frac{3}{\sqrt{2}} = 4\sqrt{2}
\]
\[
EM_{\min} = EO - R = \frac{5}{\sqrt{2}} - \frac{3}{\sqrt{2}} = \sqrt{2}
\]
故 $EM \in [\sqrt{2}, 4\sqrt{2}]$。

\textbf{（2）求二面角 $A-l-C$ 的最小值：}\\
因为 $AE \parallel CD$，$AE \subset$ 平面 $PAE$，$CD \subset$ 平面 $PCD$，两平面的交线为 $l$，所以 $l \parallel AE \parallel CD$。\\
因为 $AC \perp AE$，所以垂直于 $l$ 的截面可以看作由 $AC$ 方向与垂直于平面 $ACDE$ 的法线方向组成。\\
设 $P$ 到平面 $ACDE$ 的距离为 $h$。设 $G$ 为 $AC$ 的中点，因为 $\triangle PAC$ 是边长为 $2\sqrt{3}$ 的正三角形，所以 $PG = 3$（$P$ 到 $AC$ 的距离恒为 $3$）。在翻折过程中，$0 \le h \le 3$。\\
在垂直于 $l$ 的截面内，平面 $PAE$ 与平面 $PCD$ 的截线方向向量可分别等效为 $\vec{n}_1 = (\sqrt{3}, h)$ 和 $\vec{n}_2 = (-\sqrt{3}, h)$（以 $G$ 为基准，分别向 $A,C$ 延伸 $\sqrt{3}$ 的水平距离，并有 $h$ 的垂直高度）。\\
设二面角为 $\omega$，则：
\[
\cos\omega = \frac{\vec{n}_1 \cdot \vec{n}_2}{|\vec{n}_1||\vec{n}_2|} = \frac{-\sqrt{3} \cdot \sqrt{3} + h \cdot h}{\sqrt{3+h^2}\sqrt{3+h^2}} = \frac{h^2 - 3}{h^2 + 3}
\]
函数 $f(h) = \frac{h^2 - 3}{h^2 + 3} = 1 - \frac{6}{h^2 + 3}$ 在 $h \in [0, 3]$ 上随 $h$ 增大而增大。\\
所以当 $h$ 取得最大值 $h=3$（即平面 $PAC \perp$ 平面 $ACDE$）时，$\cos\omega$ 取得最大值 $\frac{9-3}{9+3} = \frac{1}{2}$。\\
由于 $\omega \in (0, \pi)$，$\cos\omega$ 越大，$\omega$ 越小。\\
故 $\omega_{\min} = \frac{\pi}{3}$。即二面角 $A-l-C$ 的最小值为 $\frac{\pi}{3}$。
}

\newcommand{\qSolidFoldingPentagonMinAngleTeachingNotes}{
\item \textbf{结构洞察}：本题的核心在于抓住不变量和几何结构。一是 $AE \parallel CD \implies l \parallel AE \parallel CD$ 的平行模型，从而将二面角转化为一个关于高度 $h$ 的代数最值问题。
\item \textbf{降维打击}：第一问中 $\overrightarrow{AH} \cdot \overrightarrow{CF} = 0$ 利用平面向量非常巧妙，避开了复杂的空间建系建坐标。建议在教案中引导学生观察 $F, H$ 投影到 $AC$ 上都在同一点 $G$。
\item \textbf{二面角的最值}：可以作为“建系与向量法”的进阶拔高。把截面上的方向向量抽离出来：$\vec{v}_1 = (\sqrt{3}, h)$ 和 $\vec{v}_2 = (-\sqrt{3}, h)$，利用 $\cos\omega = \frac{\vec{v}_1 \cdot \vec{v}_2}{|\vec{v}_1||\vec{v}_2|}$ 分析函数单调性。
}


\begin{document}

\begin{center}
  {\zihao{2}\bfseries 拔高专题 \quad 立体几何折叠与折痕问题专题讲义}\\[3mm]
  {\Large \bfseries 从不变量建立到空间量转化}\\[2mm]
  {\large 适用：高一/高二一轮复习 / 拔高专题课 / 120 分钟教案}\\[2mm]
  \rule{0.85\linewidth}{0.5pt}
\end{center}

\begin{teachbox}{专题设计定位}
立体几何中的“图形折叠问题”是高考与自主招生的一大高频难点。传统的授课方式往往按照“线面垂直、体积计算、点面距离、轨迹长度”进行机械平铺，导致学生在面对“动态折叠”时产生严重的畏难情绪。\\
本专题本着\textbf{“折叠 = 绕折痕旋转”}的唯一核心主线，打破常规的板块划分，以\textbf{学生认知门槛}为梯度重新编排，沿以下路径推进教学：\\
$$\text{折痕} \rightarrow \text{垂足} \rightarrow \text{定长} \rightarrow \text{截面圆弧} \rightarrow \text{空间量转化}$$
通过六层递进式的认知台阶，使学生建立起由浅入深的逻辑链条，在面对任何折叠题目时，都能利用同一套“不变量”快速破局。
\end{teachbox}

% ==================== 教学大纲 ====================
\section{120分钟教学大纲与口诀}

\begin{center}
\renewcommand{\arraystretch}{1.2}
\begin{tabularx}{0.95\linewidth}{>{\bfseries}p{2.2cm} c p{4.5cm} X}
\toprule
认知层级 & 建议时间 & 核心教学内容 & 教学目标与关键动作 \\
\midrule
核心不变量 & 15min & 折痕、垂足、定长、圆弧 & 掌握旋转模型本质，建立折前与折后守恒量。 \\
位置证明 & 20min & 垂直不变量的线面垂直转化 & 利用 $\text{折前垂直} \implies \text{折后垂直}$ 构建线面垂直。 \\
距离转化 & 25min & 中点换点、平行换点与换面 & 避开空间垂线，降维转化回平面计算距离。 \\
体积最值 & 20min & 高度最大值与垂直临界状态 & 理解体积由高决定，最值发生于折叠面垂直底面时。 \\
夹角最值 & 20min & 线面角定义与三正弦定理 & 利用 $\tan\theta = \frac{h}{d_{\text{proj}}}$ 或三正弦公式进行代数求导/二次函数最值。 \\
轨迹圆弧 & 15min & 绕垂足圆弧与弦长角度转化 & 将空间距离范围转化为圆心角，求动点轨迹弧长。 \\
综合压轴 & 选讲/挑战 & 外接球、动点最值、二面角截面 & 融合前六层模型，作为拔高检验或课后挑战题。 \\
总结复盘 & 5min & 师生共建决策图与解题口诀 & 提炼核心解题步骤。 \\
\bottomrule
\end{tabularx}
\end{center}

\begin{keybox}{解题口诀}
  \textbf{先找折痕，再找垂足；定长成圆，空间转平面；距离先换点，最值看高度，角度看投影。}
\end{keybox}

\newpage

% ==================== 第一层 ====================
\section{第一层：建立折叠不变量}

\begin{keybox}{核心思维——不变量思维：以静制动}
折叠题目的千变万化，都是由于面与面之间“角度”在发生变化。但无论角度如何变，以下五大几何量或几何关系在折叠前后\textbf{绝对保持守恒}：
\begin{enumerate}[label=\arabic*., leftmargin=2em, itemsep=0.1em]
  \item \textbf{折痕不动}：折痕所在的直线及其线段在空间中的位置与长度恒定。
  \item \textbf{折痕上的点不动}：折痕上的点是旋转的静止中心或静止轴点。
  \item \textbf{面内长度与夹角不变}：被折叠图形内部的线段长度、夹角、面积完全保持不变。
  \item \textbf{垂直折痕的线折后仍垂直}：若折叠前有直线 $l \perp$ 折痕于 $O$，则折后 $l' \perp$ 折痕于 $O$。这是确定旋转面的基础。
  \item \textbf{动点绕折痕旋转以垂足为圆心}：折起面上的点 $P$ 绕折痕旋转，其轨迹是一个圆或圆弧，圆心即为 $P$ 到折痕的垂足 $O$，半径为 $PO$。
\end{enumerate}
\end{keybox}

\par\addvspace{0.4em}{\centering
\begin{tikzpicture}[scale=0.9, >=stealth, line join=round, line cap=round]
  % 左图：折叠前
  \begin{scope}[xshift=0cm, yshift=0cm]
    \coordinate (B) at (0,0);
    \coordinate (C) at (2.5,0);
    \coordinate (D) at (2.5,2.5);
    \coordinate (A) at (0,2.5);
    \coordinate (O) at (1.25,1.25);
    
    \draw[thick] (A) -- (B) -- (C) -- (D) -- cycle;
    \draw[thick, dashed, blue] (B) -- (D);
    \draw[thick, dashed, red] (A) -- (O);
    
    \fill (A) circle (1.2pt) node[above left] {\small $A$};
    \fill (B) circle (1.2pt) node[below left] {\small $B$};
    \fill (C) circle (1.2pt) node[below right] {\small $C$};
    \fill (D) circle (1.2pt) node[above right] {\small $D$};
    \fill (O) circle (1.2pt) node[below right] {\small $O$};
    
    \draw[gray, thin] ($(O)!0.1!(A)$) -- ($($(O)!0.1!(A)$) + ($(O)!0.1!(D)$) - (O)$) -- ($(O)!0.1!(D)$);
    \node[below] at (1.25, -0.2) {\small 折叠前（正方形）};
  \end{scope}

  \draw[->, very thick, gray] (3.0, 1.25) -- (4.0, 1.25) node[midway, above] {\small 折叠};

  % 右图：折叠后
  \begin{scope}[xshift=4.8cm, yshift=0cm]
    \coordinate (B) at (0,0);
    \coordinate (C) at (2.5,0);
    \coordinate (D) at (2.9,0.8);
    \coordinate (O) at (1.45,0.4);
    \coordinate (P) at (1.1,1.5); % 旋转折起后的 A 点，记为 P
    
    \draw[thick] (B) -- (C) -- (D);
    \draw[thick, dashed, gray!60] (D) -- (B);
    
    \draw[thick, blue] (P) -- (B);
    \draw[thick, blue] (P) -- (D);
    \draw[thick, dashed, red] (P) -- (O);
    \draw[thick, dashed, gray!60] (O) -- (C);
    
    \fill (B) circle (1.2pt) node[below left] {\small $B$};
    \fill (C) circle (1.2pt) node[below right] {\small $C$};
    \fill (D) circle (1.2pt) node[above right] {\small $D$};
    \fill (O) circle (1.2pt) node[below] {\small $O$};
    \fill (P) circle (1.2pt) node[above] {\small $P\ (A')$};
    
    % 画旋转圆弧轨迹
    \draw[dashed, red, thin] (O) ++(60:1.1) arc (60:120:1.1);
    \draw[red, thin] ($(O)!0.08!(P)$) -- ($($(O)!0.08!(P)$) + ($(O)!0.08!(D)$) - (O)$) -- ($(O)!0.08!(D)$);
    
    \node[below] at (1.45, -0.2) {\small 折叠后（$PO = AO, PO \perp BD$）};
  \end{scope}
\end{tikzpicture}\par}\addvspace{0.4em}

\begin{exbox}{例 1.1}
在边长为 $2$ 的正方形 $ABCD$ 中，对角线 $AC$ 与 $BD$ 相交于点 $O$，将 $\triangle ABD$ 沿对角线 $BD$ 折起，使得点 $A$ 到达点 $P$ 的位置。
\begin{enumerate}[label=(\arabic*), leftmargin=2em, itemsep=0.1em]
  \item 证明：无论折起何种角度，均有 $PO \perp BD$；
  \item 若折起后 $PC = \sqrt{2}$，求此时线段 $PO$ 的长度以及 $\triangle PBD$ 的面积。
\end{enumerate}
\end{exbox}

\begin{paracol}{2}
\teachblock{标准解答与讲解主线}{
  \textbf{（1）证明：}\\
  In折叠前，因为四边形 $ABCD$ 为正方形，所以对角线 $AC \perp BD$。\\
  因为点 $O$ 是对角线 $AC$ 与 $BD$ 的交点，所以有 $AO \perp BD$。\\
  当将 $\triangle ABD$ 沿折痕 $BD$ 折起后，由于 $\triangle ABD$ 在旋转过程中其形状与大小保持不变，即 $\triangle PBD \cong \triangle ABD$。\\
  因此，折叠前的垂直关系在折叠后依然保持不变，即：
  \[ \boxed{PO \perp BD} \]
}
\switchcolumn
\teachblock{易错分析与教学提示}{
  \textbf{学生常见误区：}\\
  初学者往往在图形折起后，直观感觉空间中的某些线段“缩短”了或“变长”了，从而误认为折叠面内部的长度发生了变化。
}
\switchcolumn*
\teachblock{标准解答与讲解主线 (续)}{
  \textbf{（2）解：}\\
  因为正方形的边长为 $2$，其对角线长度为：
  \[ AC = BD = 2\sqrt{2} \]
  由于对角线互相垂直平分，所以：
  \[ AO = CO = \frac{1}{2} AC = \sqrt{2} \]
  由于折叠不变量性质，线段长度保持不变，故折起后有：
  \[ \boxed{PO = AO = \sqrt{2}} \]
  在三角形 $\triangle PBD$ 中，底边为折痕 $BD = 2\sqrt{2}$，由（1）可知高线为 $PO = \sqrt{2}$。\\
  故 $\triangle PBD$ 的面积为：
  \[ S_{\triangle PBD} = \frac{1}{2} \cdot BD \cdot PO = \frac{1}{2} \cdot 2\sqrt{2} \cdot \sqrt{2} = \boxed{2} \]
}
\switchcolumn
\teachblock{易错分析与教学提示 (续)}{
  \textbf{启发与提问口径：}\\
  \begin{itemize}[leftmargin=1.2em, itemsep=0.3em, topsep=0.2em]
    \item “当我们将这半张纸折起来时，这半张纸本身有没有被扯大或撕裂？没有。它的形状没有任何变化。所以只要是原本就在这半张纸上的线段长度和夹角（如 $PO$ 和 $\angle PBD$），折叠后完全保持原样。”
    \item “折叠不变量是所有折叠题的‘根’，第一步一定要在图形上把这些不变的长度 and 垂直关系用彩笔标出来。”
  \end{itemize}
}
\end{paracol}

\newpage

% ==================== 第二层 ====================
\section{第二层：位置关系证明}

\begin{keybox}{核心思维——垂直转化：从平面到空间的跨越}
折叠题目的几何证明（线面垂直与面面垂直），其核心入口非常固定，即\textbf{“利用折痕两边的垂直不变量构建线面垂直”}。
标准证明链条如下：
\[
\begin{aligned}
\text{折叠前：} \quad & AO \perp \text{折痕 } BD \\
\text{折叠后：} \quad & PO \perp \text{折痕 } BD \quad (\text{不变量}) \\
\text{底面中：} \quad & CO \perp \text{折痕 } BD \quad (\text{底面原有}) \\
& PO \cap CO = O \quad \Longrightarrow \quad \boxed{BD \perp \text{平面 } POC}
\end{aligned}
\]
这一法则适合于几乎所有的折叠垂直证明。教学重点是让学生明白：\textbf{“寻找垂直于折痕的两条相交线，是解题最天然的突破口”}。
\end{keybox}

\par\addvspace{0.4em}{\centering
\begin{tikzpicture}[scale=0.9, >=stealth, line join=round, line cap=round]
  \coordinate (B) at (0,0);
  \coordinate (C) at (2.5,0);
  \coordinate (D) at (2.9,0.8);
  \coordinate (O) at (1.45,0.4);
  \coordinate (P) at (1.1,1.5); 
  
  \draw[thick] (B) -- (C) -- (D);
  \draw[thick, dashed, gray!60] (D) -- (B);
  
  \draw[thick, blue] (P) -- (B);
  \draw[thick, blue] (P) -- (D);
  \draw[thick, dashed, red] (P) -- (O);
  \draw[thick, dashed, red] (O) -- (C);
  \draw[thick, dashed, red] (P) -- (C);
  
  \fill (B) circle (1.2pt) node[below left] {\small $B$};
  \fill (C) circle (1.2pt) node[below right] {\small $C$};
  \fill (D) circle (1.2pt) node[above right] {\small $D$};
  \fill (O) circle (1.2pt) node[below] {\small $O$};
  \fill (P) circle (1.2pt) node[above] {\small $P$};
  
  % 直角标记
  \draw[red, thin] ($(O)!0.08!(P)$) -- ($($(O)!0.08!(P)$) + ($(O)!0.08!(D)$) - (O)$) -- ($(O)!0.08!(D)$);
  \draw[red, thin] ($(O)!0.08!(C)$) -- ($($(O)!0.08!(C)$) + ($(O)!0.08!(D)$) - (O)$) -- ($(O)!0.08!(D)$);
  
  \node[below] at (1.45, -0.2) {\small $BD \perp PO$ 且 $BD \perp CO \implies BD \perp$ 平面 $POC$};
\end{tikzpicture}\par}\addvspace{0.4em}

\begin{exbox}{例 2.1}
在等腰直角三角形 $ABC$ 中，$\angle A = 90^\circ$，$AB = AC = 2$，$D, E$ 分别为 $AB, AC$ 的中点。将 $\triangle ADE$ 沿折痕 $DE$ 折起，使得点 $A$ 到达点 $P$ 的位置，且点 $P$ 在底面 $BCED$ 上的射影 $H$ 恰好是线段 $BC$ 的中点。记线段 $DE$ 的中点为 $M$。
\begin{enumerate}[label=(\arabic*), leftmargin=2em, itemsep=0.1em]
  \item 证明：$DE \perp$ 平面 $PMH$；
  \item 证明：平面 $PBC \perp$ 平面 $PMH$。
\end{enumerate}
\end{exbox}

\begin{paracol}{2}
\teachblock{标准解答与讲解主线}{
  \textbf{（1）证明：}\\
  在折叠前，因为 $\triangle ABC$ 为等腰直角三角形，且 $AB=AC$，点 $D, E$ 分别为 $AB, AC$ 的中点，\\
  所以 $DE \parallel BC$，且 $\triangle ADE$ 也是等腰直角三角形。\\
  因为点 $M$ 是 $DE$ 的中点，根据对称性，有 $AM \perp DE$。\\
  当折起后点 $A$ 到达点 $P$ 的位置，由折叠不变量性质可知：
  \[ PM \perp DE \quad \text{—— 关系 (1)} \]
  因为点 $P$ 在底面 $BCED$ 上的射影为 $H$，所以有：
  \[ PH \perp \text{平面 } BCED \]
  因为直线 $DE \subset$ 平面 $BCED$，所以有：
  \[ PH \perp DE \quad \text{—— 关系 (2)} \]
  在平面 $PMH$ 内，$PM \cap PH = P$，且 $PM, PH$ 是两条相交直线。\\
  结合关系 (1) 与 (2)，由线面垂直判定定理，可知：
  \[ \boxed{DE \perp \text{平面 } PMH} \]
}
\switchcolumn
\teachblock{易错分析与教学提示}{
  \textbf{学生常见误区：}\\
  证明（1）时，学生容易漏写相交条件 $PM \cap PH = P$。在证明（2）时，容易分不清是“线垂直面”推出“面垂直面”，还是直接用夹角证，写得不够严谨。
}
\switchcolumn*
\teachblock{标准解答与讲解主线 (续)}{
  \textbf{（2）证明：}\\
  在折叠前，因为 $DE \parallel BC$，所以折叠后在底面内仍有：
  \[ DE \parallel BC \]
  由（1）中已证：$DE \perp$ 平面 $PMH$。\\
  因为 $DE \parallel BC$，根据平行线垂直于同一平面的性质，可得：
  \[ BC \perp \text{平面 } PMH \]
  因为直线 $BC \subset$ 平面 $PBC$，由面面垂直判定定理，可知：
  \[ \boxed{\text{平面 } PBC \perp \text{平面 } PMH} \]
}
\switchcolumn
\teachblock{易错分析与教学提示 (续)}{
  \textbf{启发与提问口径：}\\
  \begin{itemize}[leftmargin=1.2em, itemsep=0.3em, topsep=0.2em]
    \item “我们想证 $DE \perp$ 平面 $PMH$，需要在平面内找两条和 $DE$ 垂直的相交线。第一条是折叠面里的 $PM$（这是折前就垂直的不变量）；第二条是高度线 $PH$（因为它垂直于整个底面，所以垂直于底面内一切线）。”
    \item “只要把‘折前垂直’与‘高线垂直’联立起来，线面垂直就必然成立。这是折叠证明题雷打不动的标准套路。”
  \end{itemize}
}
\end{paracol}

\newpage

% ==================== 第三层 ====================
\section{第三层：距离题}

\begin{keybox}{核心思维——距离转换：化难为易，降维求精}
在折起后的三维多面体中，要求某个点到平面的距离 $d$，如果直接作空间垂线，垂足往往飘落在空中或某条线上，极难证明和定位。
我们的通用战术是\textbf{“转换点与面，回原平面计算”}：
\begin{enumerate}[label=\arabic*. , leftmargin=2em, itemsep=0.1em]
  \item \textbf{利用中点/比例转换点}：若 $M$ 是 $PC$ 的中点，则 $d(M, \text{面}) = \frac{1}{2} d(C, \text{面})$。
  \item \textbf{利用平行转换点}：若直线 $l \parallel$ 平面 $\alpha$，则 $l$ 上任意点到平面 $\alpha$ 的距离都相等。
  \item \textbf{利用等体积换底法}：$V = \frac{1}{3} S_1 h_1 = \frac{1}{3} S_2 h_2$，避免直接作高。
  \item \textbf{利用面面垂直降维}：若平面 $\alpha \perp$ 平面 $\beta$，交线为 $l$，则点 $P \in \alpha$ 到平面 $\beta$ 的距离，直接等于点 $P$ 在平面 $\alpha$ 内到交线 $l$ 的距离。
\end{enumerate}
\end{keybox}

\begin{exbox}{例 3.1}
在边长为 $2$ 的正方形 $ABCD$ 中，点 $O$ 为对角线 $AC$ 与 $BD$ 的交点。现将 $\triangle ABD$ 沿对角线 $BD$ 折起，使得点 $A$ 到达点 $P$ 的位置，且满足 $PC = 2$。若点 $M$ 为线段 $PC$ 的中点，求点 $M$ 到平面 $PBD$ 的距离。
\end{exbox}

\begin{paracol}{2}
\teachblock{标准解答与讲解主线}{
  \textbf{步骤一：转换目标点}\\
  因为点 $M$ 是线段 $PC$ 的中点，且直线 $PC$ 与平面 $PBD$ 相交于点 $P$。\\
  根据平行投影与比例关系，点 $M$ 到平面 $PBD$ 的距离等于点 $C$ 到平面 $PBD$ 距离的一半：
  \[ d(M, \text{平面 } PBD) = \frac{1}{2} d(C, \text{平面 } PBD) \quad \text{—— 关系 (1)} \]
  
  \textbf{步骤二：寻找并利用面面垂直降维}\\
  在折叠前，正方形的对角线互相垂直平分，故有 $AO \perp BD$ 且 $CO \perp BD$。\\
  折叠后，点 $A$ 到达点 $P$ 处。由折叠不变量，有：
  \[ PO \perp BD \quad \text{且} \quad CO \perp BD \]
  又 $PO \cap CO = O$，故有 $BD \perp$ 平面 $POC$。\\
  因为直线 $BD \subset$ 平面 $PBD$，由面面垂直判定定理，可知：
  \[ \text{平面 } PBD \perp \text{平面 } POC \]
  两平面的交线为 $PO$。\\
  根据面面垂直的性质定理，点 $C$ 到平面 $PBD$ 的距离，即为在平面 $POC$ 内，点 $C$ 到交线 $PO$ 的距离。
}
\switchcolumn
\teachblock{易错分析与教学提示}{
  \textbf{学生常见误区：}\\
  学生常常倾向于“哪里有问，就在哪里作图”，直接在空间中去作 $M$ 到平面 $PBD$ 的垂线，导致画出很多多余的干扰线，最后无法进行定量证明。
}
\switchcolumn*
\teachblock{标准解答与讲解主线 (续)}{
  \textbf{步骤三：在平面三角形 $\triangle POC$ 内计算}\\
  在正方形 $ABCD$ 中，$AO = CO = \sqrt{2}$，折起后有 $PO = AO = \sqrt{2}$。\\
  在 $\triangle POC$ 中，已知：
  \[ PO = \sqrt{2},\quad CO = \sqrt{2},\quad PC = 2 \]
  因为 $PO^2 + CO^2 = (\sqrt{2})^2 + (\sqrt{2})^2 = 4 = PC^2$，\\
  根据勾股定理逆定理，可知 $\triangle POC$ 是以 $O$ 为直角的直角三角形。\\
  因此，在平面 $POC$ 内，点 $C$ 到直角边 $PO$ 的距离，就等于直角边 $CO$ 的长度，即：
  \[ d(C, \text{平面 } PBD) = CO = \sqrt{2} \]
  
  \textbf{步骤四：代回求最终解}\\
  将此距离代回关系 (1) 中，可得点 $M$ 到平面 $PBD$ 的距离为：
  \[ d(M, \text{平面 } PBD) = \frac{1}{2} \cdot \sqrt{2} = \boxed{\frac{\sqrt{2}}{2}} \]
}
\switchcolumn
\teachblock{易错分析与教学提示 (续)}{
  \textbf{启发与提问口径：}\\
  \begin{itemize}[leftmargin=1.2em, itemsep=0.3em, topsep=0.2em]
    \item “我们能直接作 $M$ 的垂线吗？在空间里很难确定它落在哪。但是我们发现 $M$ 是 $PC$ 的中点，根据相似比，只要算出端点 $C$ 到平面的距离，除以 2 就可以了。这就是‘换点法’。”
    \item “要找 $C$ 到平面 $PBD$ 的距离，正好发现对角线截面 $POC$ 垂直于平面 $PBD$。这样我们就把一个‘空间点到面的距离’降维变成了‘平面直角三角形内点到线的距离’，也就是 $CO$ 的长度。计算瞬间就化简了。”
  \end{itemize}
}
\end{paracol}

\newpage

% ==================== 第四层 ====================
\section{第四层：体积与表面积最值}

\begin{keybox}{核心思维——高度最值与垂直临界}
三棱锥（四面体）的体积计算公式为：
\[ V = \frac{1}{3} S_{\text{底}} \cdot h \]
当底面在折叠过程中形状与面积保持固定时，体积的最大值完全由\textbf{折叠点到底面的高度 $h$} 决定。\\
由于动点绕折痕做圆周运动，其高度 $h$ 不超过旋转的半径（即折前顶点到折痕的距离）。\\
因此：
\begin{enumerate}[label=\arabic*. , leftmargin=2em, itemsep=0.1em]
  \item \textbf{高度最大状态}：当折起面与底面所成的二面角为 $90^\circ$（即\textbf{折起面 $\perp$ 底面}）时，高度 $h$ 达到最大值，此时体积最大。
  \item \textbf{注意区分}：必须要向学生反复强调，\textbf{“折起面与底面垂直”只是体积最大时的特定最值状态，而不是折起后的普遍状态}。很多同学会在非最值的常规问题中，也想当然地套用面面垂直，导致严重失分。
\end{enumerate}
\end{keybox}

\begin{exbox}{例 4.1}
在等腰直角三角形 $ABC$ 中，$\angle C = 90^\circ$，$AC = BC = 2$，点 $D$ 是斜边 $AB$ 上的一点且满足 $AD = 1$。过点 $D$ 作 $DE \perp AC$ 于点 $E$。现将 $\triangle ADE$ 沿折痕 $DE$ 折起，使得点 $A$ 到达点 $P$ 的位置，构成四棱锥 $P-BCED$。求四棱锥 $P-BCED$ 体积的最大值。
\end{exbox}

\begin{paracol}{2}
\teachblock{标准解答与讲解主线}{
  \textbf{步骤一：求底面梯形 $BCED$ 的面积}\\
  在等腰 Rt$\triangle ABC$ 中，$AC=BC=2$，故 $\angle A = 45^\circ$。\\
  因为 $DE \perp AC$，所以在 Rt$\triangle ADE$ 中，$\angle ADE = 45^\circ$，有 $AE = DE$。\\
  因为斜边段 $AD = 1$，所以：
  \[ AE = DE = AD \cos 45^\circ = \frac{\sqrt{2}}{2} \]
  则底面梯形 $BCED$ 的高为：
  \[ CE = AC - AE = 2 - \frac{\sqrt{2}}{2} \]
  底面梯形 $BCED$ 的两底边为 $DE = \frac{\sqrt{2}}{2}$ 且 $BC = 2$。\\
  故底面面积为：
  \[ S_{BCED} = \frac{1}{2}(DE + BC) \cdot CE = \frac{1}{2}\left(\frac{\sqrt{2}}{2} + 2\right)\left(2 - \frac{\sqrt{2}}{2}\right) \]
  利用平方差公式：
  \[ S_{BCED} = \frac{1}{2}\left(2^2 - \left(\frac{\sqrt{2}}{2}\right)^2\right) = \frac{1}{2}\left(4 - \frac{1}{2}\right) = \frac{7}{4} \]
}
\switchcolumn
\teachblock{易错分析与教学提示}{
  \textbf{学生常见误区：}\\
  学生容易算错底面面积。另外，部分学生不能理解“折叠面与底面垂直时，高最大”的物理直观。
}
\switchcolumn*
\teachblock{标准解答与讲解主线 (续)}{
  \textbf{步骤二：分析高度最值}\\
  折起后，四棱锥 $P-BCED$ 的体积为：
  \[ V = \frac{1}{3} S_{BCED} \cdot h = \frac{7}{12} h \]
  其中 $h$ 是顶点 $P$ 到平面 $BCED$ 的距离。\\
  因为折叠过程是 $\triangle ADE$ 绕折痕 $DE$ 旋转，点 $P$ 的轨迹是以 $E$ 为圆心、$AE$ 为半径的圆弧。\\
  因此，点 $P$ 到平面 $BCED$ 的最大高度 $h$ 不能超过旋转半径：
  \[ h \le AE = \frac{\sqrt{2}}{2} \]
  当且仅当平面 $PDE \perp$ 平面 $BCED$ 时，高度达到最大值 $h_{\max} = \frac{\sqrt{2}}{2}$。
  
  \textbf{步骤三：计算最大体积}\\
  此时体积取得最大值：
  \[ V_{\max} = \frac{7}{12} \cdot h_{\max} = \frac{7}{12} \cdot \frac{\sqrt{2}}{2} = \boxed{\frac{7\sqrt{2}}{24}} \]
}
\switchcolumn
\teachblock{易错分析与教学提示 (续)}{
  \textbf{启发与提问口径：}\\
  \begin{itemize}[leftmargin=1.2em, itemsep=0.3em, topsep=0.2em]
    \item “点 $P$ 绕着折痕 $DE$ 旋转，它在空间中的高度是不断变化的。就像荡秋千一样，秋千荡到最高点时（也就是秋千架所在平面和地面垂直时），高度才最大。所以 $P$ 到底面的高度最大就是旋转半径 $AE$。此时折叠面和底面恰好垂直。”
    \item “这一最值状态在选择填空题中是秒杀体积最值的核心技巧。”
  \end{itemize}
}
\end{paracol}

\newpage

% ==================== 第五层 ====================
\section{第五层：夹角最值}

\begin{keybox}{核心思维——夹角投影与三正弦定理}
求折叠中的线面角或二面角最值，通常有两大通法：
\begin{enumerate}[label=\arabic*., leftmargin=2em, itemsep=0.1em]
  \item \textbf{正切定义法}：$\tan\theta = \frac{h}{d_{\text{proj}}}$。通过建立关于折起参数（如线段长或旋转角）的代数式，将其转化为求代数式的最值问题。
  \item \textbf{三正弦定理法（跨面投影模型）}：
    \[ \sin(\text{线面角}) = \sin(\text{二面角}) \cdot \sin(\text{线线角}) \]
    即：$\sin\alpha = \sin\beta \cdot \sin\theta$。\\
    若直线在底面内（如 $CD$），且折叠面绕折痕 $BM$ 旋转，二面角为 $\beta$，直线与折痕的夹角为 $\theta$。利用此公式可以\textbf{瞬间避开空间垂线与投影的定位}，直接转化为代数形式，极大降低作图难度。
\end{enumerate}
\end{keybox}

\begin{exbox}{例 5.1（山东省实验中学 2025 高一月考真题）}
  \qSolidFoldingRectangleStem
  \par\addvspace{0.4em}{\centering
  \qSolidFoldingRectangleFigure[0.9]
  \par}\addvspace{0.4em}
\end{exbox}

\begin{paracol}{2}
\teachblock{标准解答与讲解主线}{
  \qSolidFoldingRectangleSolution
}
\switchcolumn
\teachblock{易错分析与教学提示}{
  \begin{itemize}[leftmargin=1.5em]
    \qSolidFoldingRectangleTeachingNotes
  \end{itemize}
}
\end{paracol}

\newpage

% ==================== 第六层 ====================
\section{第六层：动点轨迹}

\begin{keybox}{核心思维——轨迹降维与圆心角转换}
折叠题目的轨迹问题，其物理本质是\textbf{“动点绕折痕上的垂足做圆周运动，其轨迹为空间中的圆弧”}。
求解轨迹弧长的核心五步环流如下：
\[
\boxed{\text{旋转半径 } R} 
\rightarrow 
\boxed{\text{空间距离约束}} 
\rightarrow 
\boxed{\text{余弦定理列式}} 
\rightarrow 
\boxed{\text{旋转角 } \phi \text{ 范围}} 
\rightarrow 
\boxed{\text{弧长 } l = R\phi}
\]
解题关键是利用余弦定理，把空间中两个动点的距离（皮）与旋转平面角（骨）连接起来。
\end{keybox}

\begin{exbox}{例 6.1}
在边长为 $2$ 的正方形 $ABCD$ 中，点 $O$ 为对角线 $AC$ 与 $BD$ 的交点。现将 $\triangle ABD$ 沿对角线 $BD$ 折起，使得点 $A$ 到达点 $P$ 的位置。在折起过程中，当线段 $PC$ 的长度在 $[\sqrt{2}, \sqrt{6}]$ 范围内时，求点 $P$ 运动的轨迹对应的圆弧长度。
\end{exbox}

\begin{paracol}{2}
\teachblock{标准解答与讲解主线}{
  \textbf{步骤一：确定旋转圆心与半径}\\
  在折叠前，正方形 $ABCD$ 的对角线互相垂直平分，故 $AO \perp BD$ 且 $CO \perp BD$。\\
  折叠后，点 $A$ 绕折痕 $BD$ 旋转，旋转圆心为垂足 $O$，旋转半径为：
  \[ R = PO = AO = \sqrt{2} \]
  并且 $CO = \sqrt{2}$ 保持恒定。
  
  \textbf{步骤二：用余弦定理建立空间距离与旋转角的关系}\\
  因为 $PO \perp BD$ 且 $CO \perp BD$，所以二面角 $P-BD-C$ 的平面角为 $\phi = \angle POC$。\\
  在 $\triangle POC$ 中，由余弦定理有：
  \[ PC^2 = PO^2 + CO^2 - 2 \cdot PO \cdot CO \cos\phi \]
  代入已知长度 $PO = \sqrt{2}, CO = \sqrt{2}$ 得：
  \[ PC^2 = 2 + 2 - 2 \cdot \sqrt{2} \cdot \sqrt{2} \cos\phi = 4 - 4\cos\phi \]
}
\switchcolumn
\teachblock{易错分析与教学提示}{
  \textbf{学生常见误区：}\\
  学生看到“空间线段 $PC$ 长度范围”时，通常感到无从下手，无法将其转化为角度。他们容易忘记在 $\triangle POC$ 中直接使用余弦定理进行搭桥。
}
\switchcolumn*
\teachblock{标准解答与讲解主线 (续)}{
  \textbf{步骤三：求旋转角 $\phi$ 的变化范围}\\
  已知 $PC \in [\sqrt{2}, \sqrt{6}]$，所以 $PC^2 \in [2, 6]$。\\
  列出不等式组：
  \[ 2 \le 4 - 4\cos\phi \le 6 \implies -2 \le -4\cos\phi \le 2 \]
  同除以 $-4$，改变不等号方向：
  \[ -\frac{1}{2} \le \cos\phi \le \frac{1}{2} \]
  因为折角 $\phi \in [0, \pi]$，由余弦函数的单调性，可求得旋转角 $\phi$ 的范围为：
  \[ \phi \in \left[\frac{\pi}{3}, \frac{2\pi}{3}\right] \]
  
  \textbf{步骤四：计算圆心角的改变量与弧长}\\
  旋转对应的圆心角变化量为：
  \[ \Delta\phi = \frac{2\pi}{3} - \frac{\pi}{3} = \frac{\pi}{3} \]
  点 $P$ 运动的轨迹是一段圆弧，其半径为 $R = \sqrt{2}$。\\
  故所求轨迹弧长为：
  \[ l = R \cdot \Delta\phi = \sqrt{2} \cdot \frac{\pi}{3} = \boxed{\frac{\sqrt{2}\pi}{3}} \]
}
\switchcolumn
\teachblock{易错分析与教学提示 (续)}{
  \textbf{启发与提问口径：}\\
  \begin{itemize}[leftmargin=1.2em, itemsep=0.3em, topsep=0.2em]
    \item “点 $P$ 在天空中绕着 $O$ 旋转，它跟地面上的固定点 $C$ 距离在变化。这两个动点和旋转圆心 $O$ 组成了一个三角形 $\triangle POC$。这个三角形里哪些边是不变的？$PO$ 和 $CO$ 是定长。所以变化的只有夹角 $\phi$。我们用余弦定理就能直接把 $PC$ 的长度变化和夹角变化连起来，求出角度的范围，进而求出圆弧长。”
  \end{itemize}
}
\end{paracol}

\newpage

% ==================== 第七层 ====================
\section{第七层：综合压轴模型——非标准底面折叠}

\begin{keybox}{核心铺垫——三大进阶模型}
要解决非标准折叠的综合压轴题，必须熟练掌握以下三个“桥梁”模型：
\begin{enumerate}[label=\arabic*., leftmargin=2em, itemsep=0.1em]
  \item \textbf{外接球模型}：直角三角形的外心在斜边中点。若顶点 $P$ 到底面的垂线为 $PH$，则三棱锥的外接球球心 $O$ 必定落在 $PH$ 所在直线上。
  \item \textbf{球面动点距离范围}：若点 $M$ 在以 $O$ 为球心、$R$ 为半径的球面上运动，空间中一定点 $E$ 到球面上动点 $M$ 的距离取值范围为：$EM \in [|EO-R|, EO+R]$。这是空间版的“三角不等式”。
  \item \textbf{交线方向与二面角截面}：若 $a \parallel b$ 且 $a \subset \alpha, b \subset \beta$，平面 $\alpha \cap \beta = l$，则交线 $l \parallel a \parallel b$。利用这组平行关系，作垂直于交线 $l$ 的截面，可以将动态二面角转化为截面方向向量的平面夹角。
\end{enumerate}
\end{keybox}

\begin{exbox}{例 7.1：直角梯形底面 + 正三角形翻折 + 外接球与二面角最值}
  \qSolidFoldingPentagonMinAngleStem
  \par\addvspace{0.4em}{\centering
  \qSolidFoldingPentagonMinAngleFigure[0.9]
  \par}\addvspace{0.4em}
\end{exbox}

\begin{paracol}{2}
\teachblock{标准解答与讲解主线}{
  \qSolidFoldingPentagonMinAngleSolution
}
\switchcolumn
\teachblock{易错分析与教学提示}{
  \begin{itemize}[leftmargin=1.5em]
    \qSolidFoldingPentagonMinAngleTeachingNotes
  \end{itemize}
}
\end{paracol}

\newpage

% ==================== 决策图 ====================
\section{总结复盘与解题决策树}

\par\addvspace{0.8em}{\centering
\begin{tikzpicture}[scale=0.9, >=stealth, 
    node/.style={rectangle, draw, thick, fill=white, text width=5cm, align=center, minimum height=1cm, rounded corners=3pt},
    leaf/.style={rectangle, draw, thick, fill=gray!10, text width=4.5cm, align=center, minimum height=0.8cm}]
    
    % 主干
    \node[node] (M1) at (0, 0) {\bfseries 【第一步】寻找折痕与垂足 \\ 确定不变量：折痕 $BD$, 垂足 $O$, 定长 $PO=AO$};
    
    % 分支：证明
    \node[node] (D1) at (-4, -2.5) {\bfseries 【方向 A】证明位置垂直 \\ (线面垂直 / 面面垂直)};
    \node[leaf] (L1) at (-4, -5.0) {利用不变量垂直 $PO \perp BD$ 结合底面垂直 $CO \perp BD$ 证明 $BD \perp$ 平面 $POC$。};
    
    % 分支：距离
    \node[node] (D2) at (2, -2.5) {\bfseries 【方向 B】求解空间距离 \\ (点到平面距离 $d$)};
    \node[leaf] (L2) at (2, -5.0) {中点比例换点、平行线换点；利用 $PBD \perp POC$ 将空间高度降维转化为平面直角三角形高。};
    
    % 分支：最值与轨迹
    \node[node] (D3) at (8, -2.5) {\bfseries 【方向 C】最值与轨迹问题 \\ (体积最大 / 夹角最值 / 弧长)};
    \node[leaf] (L3) at (8, -5.0) {体积最大时面面垂直（高最大）；夹角最值用三正弦或正切；轨迹在截面用余弦定理转为圆心角。};
    
    % 连线
    \draw[->, thick] (M1) -| (D1);
    \draw[->, thick] (M1) -- (D2);
    \draw[->, thick] (M1) -| (D3);
    
    \draw[->, thick] (D1) -- (L1);
    \draw[->, thick] (D2) -- (L2);
    \draw[->, thick] (D3) -- (L3);
\end{tikzpicture}\par}\addvspace{0.8em}

\begin{keybox}{解题终极思维}
  \begin{itemize}[leftmargin=1.5em]
    \item \textbf{先看折痕}：抓静止轴。
    \item \textbf{再作垂线}：定旋转圆心 $O$ 和半径 $R$。
    \item \textbf{降维平面}：所有距离与角度，一定要找对角面（如 $\triangle POC$）或旋转截面，把三维的旋转动点问题降维拉平到平面内求解。
  \end{itemize}
\end{keybox}

\end{document}
